\section{Introduction}\label{sec:introduction}

For $v>1$, let
\[
 \beta(v)=\sum_{j=0}^{\infty}\frac{(-1)^j}{(2j+1)^v}
         =1-\frac1{3^v}+\frac1{5^v}-\cdots
\]
be the Dirichlet beta function. Its values at positive odd integers
are nonzero rational multiples of powers of $\pi$. The arithmetic
of its even values is much less understood. We study how many
irrational even values must occur in an interval whose starting
index is allowed to vary.

\begin{theorem}\label{thm:one}
For every positive odd integer $m$, at least one of
\[
 \beta(m+1),\beta(m+3),\ldots,\beta(19m-1)
\]
is irrational.
\end{theorem}

\noindent\begin{minipage}{\linewidth}
An earlier version of this paper~\cite{TshiftV1} proved
Theorem~\ref{thm:one} with $67$ in place of $19$. More generally,
that paper gave a bound for any prescribed number $B$ of irrational values.
The two statements below compare that general bound with our
improved one.
In both statements, $\varphi$ is Euler's totient function and
$\log$ is the natural logarithm.

\begin{previousversion}
For every positive odd integer $m$ and every integer $B\ge1$,
at least $B$ of
\[
 \beta(m+1),\beta(m+3),\ldots,\beta\bigl(C_B^{(1)}m-1\bigr)
\]
are irrational, where
\[
 D_B^{(1)}=\sum_{d=1}^B\varphi(d),\qquad
 C_B^{(1)}=1+2\left\lceil D_B^{(1)}
                    \left(33+\frac{45}{2}\log B\right)\right\rceil.
\]
\end{previousversion}

\begin{theorem}\label{thm:intervals}
For every positive odd integer $m$ and every integer $B\ge1$,
at least $B$ of
\[
 \beta(m+1),\beta(m+3),\ldots,\beta(C_Bm-1)
\]
are irrational, where
\begin{equation}\label{eq:closed-endpoint}
 D_B=\sum_{\substack{1\le d\le2B-1\\d\text{ odd}}}\varphi(d),
 \qquad
 C_B=1+2\left\lceil D_B\left(9+\frac{40}{9}\log(2B-1)\right)\right\rceil.
\end{equation}
\end{theorem}
\end{minipage}\par\medskip

Both bounds give an upper endpoint linear in the starting
index for each count $B$. For $B=1,2,3$, the multipliers decrease
from $67,197,463$ to $19,85,229$, respectively.
The new rational functions need zeros only on grids with odd
numerators and denominators, and their poles are spaced by $2$
instead of $4$. These choices reduce the numerator degree and
the cost of clearing denominators.

Choosing parameters for particular starting indices and counts
gives the endpoints below. The uniform column uses
Theorem~\ref{thm:intervals}, which gives $19m-1$ when $B=1$.
\begin{center}
\begin{tabular}{@{}rrrr@{}}
\toprule
First even index&\multicolumn{2}{c}{Last even index}&\# irrational values\\
\cmidrule(lr){2-3}
&Specific parameters&Uniform bound&(at least)\\
\midrule
4&54&56&1\\
6&92&94&1\\
8&128&132&1\\
20&352&360&1\\
20&1344&1614&2\\
20&3538&4350&3\\
100&1838&1880&1\\
\bottomrule
\end{tabular}
\end{center}
The rational parameter choices are given in
Section~\ref{sec:endpoints}.
Appendix~\ref{app:optimization} refines the uniform one-value bound
and explains the size of its multiplier.

\subsection{Comparison with earlier results}

Let $I(s)$ count the irrational values among
$\beta(2),\beta(4),\ldots,\beta(s)$. Our earlier paper~\cite{Tcount}
proves, for every even $s\ge160$,
\begin{equation}\label{eq:earlier-count}
 I(s)\ge\left\lfloor
 \pi\sqrt{\frac{2}{135}}\sqrt{\frac{s}{\log s}}
 \right\rfloor.
\end{equation}
To obtain $B$ irrational values starting at $\beta(m+1)$ from
such a count, we must allow for all $(m-1)/2$ preceding values
to be irrational. We therefore need a count of at least
$(m-1)/2+B$. For fixed $B$, reaching this target with
\eqref{eq:earlier-count} requires an endpoint of order
$m^2\log m$: quadratic growth in the starting index, with an
additional logarithmic factor. Theorem~\ref{thm:intervals}
gives linear growth.

For example, the earlier paper also proves the sharper count
$I(3828)\ge10$. At most nine of these irrational values
can have indices below $20$, so at least one of
$\beta(20),\beta(22),\ldots,\beta(3828)$ is irrational.
The present construction reduces this upper endpoint to $352$.
The table compares endpoints guaranteeing one irrational value
from each listed starting index:
\begin{center}
\begin{tabular}{@{}rrrr@{}}
\toprule
First even index&Endpoint from~\cite{Tcount}&Endpoint from~\cite{TshiftV1}&Endpoint here\\
\midrule
4&160&198&54\\
6&364&330&92\\
20&3828&1272&352\\
100&187\,340&6632&1838\\
\bottomrule
\end{tabular}
\end{center}
The column from~\cite{Tcount} follows from the bounds
$I(160)\ge2$, $I(364)\ge3$, $I(3828)\ge10$ and
$I(187\,340)\ge50$. Discarding the first $1$, $2$, $9$ and $49$
even beta values, respectively, leaves at least one irrational
value in each interval.

The column from~\cite{TshiftV1} uses the specific parameter
choices in Section~8.2 of that paper at starting indices $4$
and $6$, and the bound $67m-1$ at $20$ and $100$.
The last column uses the individual parameter choices in
Section~\ref{sec:finite-endpoints}.

Sorokin~\cite[Theorem~S]{S07} treats shifted intervals directly.
For every sufficiently large odd $m$ and every even
$q>100m\log m$, he proves that at least one of
\[
 \beta(m+1),\beta(m+3),\ldots,\beta(m+q-1)
\]
is irrational. Huttner~\cite[Remark~5]{H12} gives the same
conclusion. Taking the smallest allowable $q$ gives an upper
endpoint of order $m\log m$. Theorem~\ref{thm:one} removes
this logarithmic factor and applies to every positive odd $m$.

At the fixed starting index $4$, Krisman~\cite{K26} gives a
shorter interval: at least one of
$\beta(4),\beta(6),\ldots,\beta(36)$ is irrational.
Earlier finite results starting at $\beta(2)$ reached the endpoint
$10$ in Lai and Zhou~\cite[Theorem~6.1]{LZ22}; those one-value
statements are now subsumed by the irrationality of
$\beta(2)$~\cite{OAI26}.

Taking $m=1$ in the present construction also improves the
finite counts in~\cite{Tcount}:
\begin{center}
\begin{tabular}{@{}rrrr@{}}
\toprule
\# irrational values&\multicolumn{3}{c}{Last even index, starting at $2$}\\
\cmidrule(lr){2-4}
(at least)&Endpoint from~\cite{Tcount}&Endpoint from~\cite{TshiftV1}&Endpoint here\\
\midrule
2&160&196&70\\
3&364&462&186\\
4&588&772&360\\
5&1030&1386&554\\
10&3828&5428&2794\\
\bottomrule
\end{tabular}
\end{center}
The column from~\cite{TshiftV1} uses that paper's
Theorem~1.1 with $m=1$.
The specific estimates in~\cite{Tcount} give smaller endpoints.
The last column uses the parameter choices in
Section~\ref{sec:finite-endpoints}.

\subsection{How the proof locates the irrational values}

We construct $B$ sequences of linear forms supported on the
desired interval, whose even upper index we denote by $s$.
The integer $m$ fixes the starting point, and $B$ fixes the count;
the sequence index $n$ tends to infinity. The forms have the shape
\[
 F_{n,b}=q_b(n)+
 \sum_{\substack{m+1\le\nu\le s\\\nu\text{ even}}}
          b^\nu\rho_\nu(n)\beta(\nu),
 \qquad b=1,3,\ldots,2B-1.
\]
Here $q_b(n)$ and $\rho_\nu(n)$ are rational.
If at most $B-1$ of the beta values in the interval were irrational,
a fixed integer combination of the $B$ forms would cancel their coefficients at
every $n$. The resulting sequence would be rational. We arrange
that clearing its denominators gives integers tending to zero,
with infinitely many nonzero terms. A nonzero integer has
absolute value at least $1$, so this is impossible.

To remove the values below $\beta(m+1)$, we differentiate $m-1$
times. For example, normalized second differentiation changes
\[
 \frac{c_i}{(t+2h)^i}
 \quad\hbox{into}\quad
 \binom{i+1}{2}\frac{c_i}{(t+2h)^{i+2}}.
\]
Reflection removes odd indices, leaving $4$ as the first even
exponent. The zeros that control the rational denominators must
survive this operation, so we start with triple zeros. In general,
we pair $(m-1)$st differentiation with zeros of multiplicity $m$.
Krisman~\cite[Sections~2--4]{K26} uses the second-derivative and
triple-zero device in the beta setting and attributes the
differentiation method to Rivoal.

Our poles are spaced by $2$. On a grid $j/b$ with $j$ and $b$
odd, shifting past the $h$th pole changes the alternating sign
by $(-1)^{bh}=(-1)^h$. This sign is the same on every grid,
so the beta coefficients retain their common row.
Zeros are needed only on these odd grids. This reduces the
numerator degree and arithmetic cost that determine the endpoints.
Integer zeros and half-integer poles, equivalent after rescaling
to this parity choice, also occur in~\cite{K26}.

The remaining task is to keep a nonzero rational constant after
elimination. At the seed index $n=1$, each choice of a polynomial
multiplier produces a vector of $B$ rational constants, one for
each grid. We construct $B$ such multipliers and put their vectors
in the columns of a matrix. A Gram representation using independent
polynomials and a positive weight proves that the matrix is
nonsingular. Combining the multipliers therefore lets us prescribe
a seed vector whose last coordinate alone is nonzero. Prime
dilation then propagates that constant to infinitely many indices.
Section~\ref{sec:worked-gram} shows the choice explicitly for two grids.

The polynomial multipliers, grouped denominator argument and
prime-dilation calculation adapt those of~\cite{Tcount}.
The rational-grid and coefficient-elimination framework appears
in Lai and Yu~\cite[Sections~2 and~5]{LY20}, who credit the earlier
elimination method of Fischler, Sprang and Zudilin, and is adapted
to real Dirichlet characters by Kumar and Tosi~\cite[Section~2]{KT26}.
The auxiliary-prime nonvanishing principle is the Archimedean
criterion of Lai and Sprang~\cite[Lemma~2.3]{LS25}.
We supply the arguments needed here.

Section~\ref{sec:construction} defines the functions and forms.
Section~\ref{sec:gram} constructs the seed, and
Sections~\ref{sec:arithmetic}--\ref{sec:dilation} establish
denominator clearing, decay and prime dilation.
Section~\ref{sec:contradiction} gives the contradiction;
Section~\ref{sec:endpoints} chooses the endpoints.
We use elementary prime valuations, standard complex analysis,
Stirling's formula and the prime number theorem.
